
\section{Related Work}
\label{sec:related}


\paragraph{Agent-based Social Simulation}
Recent work has explored LLM-powered agents for social simulation.
Generative Agents~\citep{park2023generative} pioneers this direction by simulating a society of 25 agents for two days,
observing emergent social behaviors such as organizing a party.
Humanoid Agents~\citep{wang2023humanoid} extends this by incorporating human needs inspired by Maslow's hierarchy of needs~\citep{maslow1943theory}.
Project Sid~\citep{altera2024projectsid} grounds agents in an established digital game environment, \textit{Minecraft},
showing emergent specialization, collective rule formation, and cultural transmission.
BookWorld~\citep{ran-etal-2025-bookworld} constructs multi-agent systems from fictional works,
enabling story generation through character interactions.
Aivilization~\citep{fan2026aivilization} designs a game environment to simulate agents' production and economic behaviors.
CAMEL~\citep{li2024camel} investigates role-playing multi-agent systems for collaborative problem-solving in mathematics and coding tasks.
However, existing systems such as Aivilization face limitations in enabling long-term life simulation, where most LLM calls are spent on \textit{low-level operations}, \ie, \dq{physical} interactions within the virtual environment (\eg, picking up a bottle and moving between locations), rather than long-term social interactions between agents.
Table~\ref{tab:simulation-comparison} provides a systematic comparison of these systems across key dimensions.



\begin{table*}[h]
\centering
\small
\setlength{\tabcolsep}{4pt}
\caption{Comparison of \method and existing agent society systems.
``Time Scale'' denotes the temporal scope of a simulation run.
``Env.\ FB'' indicates whether environmental feedback is rule-based or LLM-generated.
``Econ.\ Sys.'' denotes whether the system includes items and economic mechanisms.
``Career Sys.'' refers to whether agents can choose and change occupations over time.
``Group Interact'' means the system supports interactions among more than two agents.
Project Sid's numbers are from Section 4.2 of its paper.
}
\label{tab:simulation-comparison}
\resizebox{\textwidth}{!}{%
\begin{tabular}{lcccccc}
\toprule
& \textbf{Gen.\ Agents} & \textbf{Humanoid} & \textbf{Project Sid} & \textbf{BookWorld} & \textbf{Aivilization} & \textbf{Agentopia} \\
& \citep{park2023generative} & \citep{wang2023humanoid} & \citep{altera2024projectsid} & \citep{ran-etal-2025-bookworld} & \citep{fan2026aivilization} & (ours) \\
\midrule
Time Scale      & Days   & Days   & Days  & --- & Weeks  & Years  \\
Num.\ Agents       & 25     & ---    & 50 & ---   & ${\sim}$10,000     & 100    \\
Action Space    & Predefined & Predefined & Predefined & Free-form & Predefined & Free-form \\
Env.\ FB        & Rule & Rule & Rule & LLM & Rule & LLM \\
Memory          & Retrieval     & Retrieval & Retrieval & Retrieval  & Fixed Scratchpad & File-based \\
Skill Growth    & \xmark & \xmark & \xmark & \xmark & \cmark & \cmark \\
Econ.\ Sys.     & \xmark & \xmark & \cmark & \xmark & \cmark & \cmark \\
Career Sys.     & \xmark & \xmark & \cmark & \xmark & \cmark & \cmark \\
Group Interact  & \cmark & \xmark & \cmark & \cmark & \xmark & \cmark \\
Reward  & \xmark & \xmark & \xmark & \xmark & \cmark & \cmark \\
Training       & \xmark & \xmark & \xmark & \xmark & \xmark & \cmark \\
\bottomrule
\end{tabular}%
}
\end{table*}

\paragraph{Persona Simulation and LLM Role-Playing}
This line of research aims to improve LLMs' anthropomorphism and role-playing ability~\citep{chen2024from}.
\citet{shanahan2023role} argue that LLMs are inherently role-playing,
whether acting as general assistants or portraying specific characters.
\citet{wang2025coser} propose that training on diverse characters enables LLMs to generalize anthropomorphic abilities,
bridging persona simulation and role-playing research.
Prior studies focus on data, evaluation, and optimization:
For data,
Character-LLM~\citep{shao2023character} and RoleLLM~\citep{wang2023rolellm} apply LLMs to synthesize character-specific question-answer data.
CoSER~\citep{wang2025coser} extracts authentic character data from literary corpora,
obtaining large-scale, high-quality role-playing datasets.
For evaluation,
previous studies primarily employ LLM judges~\citep{wang2023rolellm, tu2024charactereval}
and introduce rubrics aligned with human preferences~\citep{wang2025coser, du2026her}.
Additionally, RoleEval~\citep{shen2023roleeval}, InCharacter~\citep{wang2024incharacter},
and LifeChoice~\citep{xu2024character}
evaluate LLMs objectively from the perspectives of knowledge, personality, and decision-making, respectively.
For optimization, 
prior research applies supervised fine-tuning with human data~\citep{wang2025coser, zhou2023characterglm}
or synthetic data~\citep{chan2024personahub, lu2024large}
CogDual~\citep{liu-etal-2025-cogdual}, CPO~\citep{ye-etal-2025-cpo}, and HER~\citep{du2026her}
explore reinforcement learning for role-playing,
where obtaining high-quality, human-aligned reward signals remains a key challenge.
Besides, Persona Vector~\citep{chen2025persona}, Assistant Axis~\citep{lu2026assistant},
and Emotion Concepts~\citep{sofroniew2026emotion}
explore the internal mechanisms of persona and emotion in LLMs,
proposing methods for detection and manipulation.


